\appendix


\subsection{Training Scale}\label{app:scale}
Holding the reward and reference fixed, we vary the number of parallel environments from $64$ to $512$ (\autoref{fig:scaling}), plotting return against total environment steps so that runs with different $N$ are compared at matched sample count. PPO minibatch size and epochs per iteration are held fixed. Final return grows with scale (dog $19\to137$, trot $17\to70$). The gain persists at matched sample count, since at equal total steps the $64$-environment runs trail the larger-$N$ runs on both clips, so the effect is not merely a batch-size artifact. Survival recorded \textit{during training} (over stochastic exploration rollouts) reaches $7.6$ s (dog) and $6.1$ s (trot) at $512$ environments, while the corresponding final policies, evaluated deterministically, survive the $19.6-20.0$ s of \autoref{tab:all}. The difference is exploration noise rather than a different policy. We vary only scale and do not claim it dominates reward shaping, which we did not sweep, and the claim is bounded to the two clips tested.

\begin{figure}[t]
\centering
\includegraphics[width=\linewidth]{figs/scaling_ablation.pdf}
\caption{Training-scale ablation. PPO return versus total environment steps for $64-512$ parallel environments on the dog-walk and horse-trot references, reward and reference held fixed. Solid lines are an exponential moving average, and the light band is the raw per-iteration return of a single run per setting.}
\label{fig:scaling}
\end{figure}

\subsection{Additional Retargeting Analyses}\label{app:retarget}
These analyses expand \autoref{subsec:exp-minimal} and \autoref{subsec:exp-abl}.

\textbf{Hip regularization.} Hip regularization is close to inert, moving the reference by $0.15^\circ$ (dog) and $0.24^\circ$ (horse). We retain it because it costs nothing and stabilizes the body-frame heading readout, rather than because the measurements demand it. The $e_q^{\mathrm{ref}}$ of the ablated dog-walk policies stays within seed noise of the full method ($13.9-15.9^\circ$ vs.\ $14.3^\circ$). On the horse trot we retrained the full, depth-without-prior and $-$hip-reg variants over three additional seeds each. The $-$hip-reg variant matches the full method ($11.7\pm0.5^\circ$ vs.\ $11.6\pm0.1^\circ$) and survives $17.5\pm3.5$ s, so the single-seed instability we had first seen there was a fluke. Depth-without-prior, by contrast, fails consistently ($29.1\pm4.5^\circ$, survival $3.0\pm2.8$ s). The knee-branch row of \autoref{fig:abl} is absent by construction, because ANYmal-D's knee sign is fixed by its kinematics and the opposite IK branch violates joint limits rather than offering an alternative retargeting.

\textbf{The depth-to-heading branch flip.} One indirect channel explains the non-zero residual of the depth-consistency check (\autoref{subsec:exp-minimal}, ``the ablation isolates one axis''). The body-frame heading is estimated during preprocessing from the four hip keypoints, before the lateral pin is applied, so zeroing the depth-derived lateral shifts that heading estimate slightly and rotates the fore-aft and vertical foot targets by a small amount. On nearly all frames this is negligible. On a few of the bow's most extreme poses the IK amplifies it into a single-joint solution-branch flip. That flip produces the $1.4^\circ$ bow RMS while leaving the median frame untouched.

\textbf{Camera roll.} Rotating the 2D keypoints in-plane simulates a camera roll of $\phi$ degrees. Re-running the pipeline drifts the reference from the $\phi{=}0$ baseline by $6.8^\circ$, $13.3^\circ$, $19.6^\circ$ and $25.3^\circ$ of joint error at $\phi=5,10,15,20^\circ$, roughly $1.3^\circ$ per degree of roll, before it falls apart at $88^\circ$ by $\phi=30^\circ$ as the corrupted foot targets force IK branch flips. The growth is smooth, but even small angles are not benign. The $13.3^\circ$ at $\phi{=}10^\circ$ already exceeds the $10.0^\circ$ RMS by which our video reference disagrees with MoCap (\autoref{subsec:exp-mocap}), so a clip with visible roll should be levelled before retargeting. This quantifies input assumption (i) of \autoref{sec:method}.

\textbf{Cost of the removed stage.} The analytic reference extraction we keep is nearly free. Preprocessing a $955$-frame clip takes about $0.2$ s, and the IK retargeting cost is dominated by a one-off simulator start-up rather than by the frames. The video-depth stage we drop is a full neural network run over every frame, at $0.075$ s per frame, roughly $70$ s for a $955$-frame clip, peaking at $3.5$ GB of GPU memory for the small model and about $7.8$ GB for the large one. Since 2D keypoint detection runs either way, removing depth saves about a minute of GPU inference and several gigabytes of memory per clip while leaving the reference essentially unchanged.

\textbf{Anatomical knee correspondence.} The anatomical ``knee'' of a quadruped is the stifle on the hind limb and the carpus/elbow complex on the fore limb, and the animal's distal metacarpal/metatarsal segment has no counterpart on ANYmal-D's three-DoF leg. Because IK is driven by paw targets alone and knee keypoints only select the solution branch (\autoref{subsec:retarget}), this extra degree of freedom is absorbed by the retargeting rather than modelled.

\textbf{Depth ablation without calibration.} The depth-without-prior variant of \autoref{fig:abl} uses raw relative disparity mapped to a root-relative coordinate up to the unknown affine scale the model does not provide, with no per-clip metric calibration, and the multi-metre lateral excursions are the direct consequence. The metric and per-clip-calibrated variants of \autoref{subsec:exp-minimal} are the fairer comparisons, and the repaired one is where the claim rests. The bias and variance argument bounds how much they can gain, since even a perfectly calibrated depth axis can at best match the constant when the true lateral variance is small.

\subsection{Additional Gait and Baseline Analyses}\label{app:baselines}
The per-behavior contact-timing numbers of \autoref{subsec:exp-sim} are collected in \autoref{tab:gait}.

\begin{table*}[t]
\centering
\scriptsize
\setlength{\tabcolsep}{4pt}
\caption{Gait-structure fidelity. Duty factor and diagonal offset for the retargeted \emph{reference} and the learned \emph{policy} in simulation, using a foot-height detector for the reference and contact force for the policy. Offset is folded to $[0,50]$, where $0$ is in-phase and $50$ antiphase.}
\label{tab:gait}
\begin{tabular}{@{}llcc@{}}
\toprule
Behavior & Quantity & Reference & Policy (sim) \\
\midrule
\multirow{2}{*}{Dog walk} & Duty factor & 0.61 & 0.73 \\
 & Diagonal offset (\% cycle) & 33 & 38 \\
\midrule
\multirow{2}{*}{Horse walk} & Duty factor & 0.61 & 0.64 \\
 & Diagonal offset (\% cycle) & 25 & 20 \todo{old appendix had 9; main text has 20} \\
\midrule
\multirow{2}{*}{Horse trot} & Duty factor & 0.57 & 0.62 \\
 & Diagonal offset (\% cycle) & 4 & 2 \\
\bottomrule
\end{tabular}
\end{table*}

\textbf{Horse-walk retraining sweep.} \todo{Reconcile before release: the previous appendix stated the horse walk's diagonal offset stayed compressed at $0-16\%$ across the weight sweep (against the reference's $25\%$), while the main text and \autoref{tab:all} report $20\pm5\%$ in simulation and $22\pm3\%$ on hardware. One of these is stale.} Retraining the horse walk across the bounded ranges of $w_r$, $w_a$ and $w_z$, a higher root weight suppresses the spurious airborne phase to $0.5\%$ but does not otherwise change the four-beat timing.

\textbf{Give-paw balance.} Balancing on three legs while one front paw swings up is only barely stable in the static sense. The center of mass briefly leaves the three-foot support triangle, so that the minimum support-polygon margin goes slightly negative and the policy holds the pose through dynamic rather than static balance. That is what makes the give-paw the least robust behavior under DR.

\textbf{Torque form of the screen.} A torque version of the actuation statistic ranks the five behaviors identically to the velocity form, with at most $2.1\%$ of samples above $80\%$ of peak torque on any behavior. We keep the velocity form in the body because it needs nothing but the reference and no rollout.

\textbf{Non-imitation velocity control (full).} To see what the imitation reward buys, we train the same setup with the imitation terms turned off (joint, foot and root weights all zero) and keep only a velocity-command tracking term set to each clip's recovered speed. This baseline keeps our observation space, action space, actuator model and DR, and retains the standard flat-terrain velocity task's own penalties on action rate and joint acceleration, but has no base-height or orientation penalty, so the crouch is a choice the reward permits rather than one it rewards. On the dog walk the resulting policy reaches the commanded speed ($0.27$ against $0.29$ m/s) and stays up the whole time, confirming that the speed, survival and airborne columns of \autoref{tab:all} can be reproduced with no imitation at all. But it gets there by crouching, dropping its base to about $0.30$ m against the reference's $0.50$, with essentially no gait fidelity ($r_j\approx0$). At the trot speed it cannot stay upright for a second.

\textbf{Naturalness (full).} The naturalness argument is measured, not asserted, by comparing the imitation policies against the non-imitation control from instrumented rollouts. The robust, seed-independent difference is posture: the velocity control reaches its speed target only by crouching, holding a mean base height of $0.28-0.30$ m against the imitation policies' upright $0.52-0.54$ m (the reference is $0.50$). On the harder horse trot the crouch is also far less efficient and far less smooth. The imitation policy has one-fifth the cost of transport ($14$ vs.\ $72$) and one-sixteenth the joint-velocity saturation ($0.5\%$ vs.\ $8.3\%$ of samples above $80\%$ of the actuator limit). On the easy dog walk the crouch is mechanically cheap and the velocity control matches or beats the imitation policy on both energy and smoothness, so those metrics do not separate them and the difference is posture alone. We therefore do not claim a uniform quantitative naturalness advantage.

\textbf{Adversarial objective (full).} AMP failed across roughly eight variants (windowed discriminators, action-rate penalties, velocity-task terms), because a style reward blind to root motion has no reason to settle on a clean gait when the references are this rough. The discriminator sees a window of $K$ consecutive frames of a $27$-dimensional feature (the twelve joint angles, the four body-frame foot positions and the body-frame root linear velocity), trained with a gradient penalty of $10$ and the same PPO backbone as our policy. The variants sweep $K\in\{2,6,12\}$, the action-rate penalty, and an added velocity-command term.

\textbf{On ground truth.} No motion capture exists for the animals in our clips, which is the premise of the work rather than a gap in it: nobody has marker-captured \textit{this} dog giving paw. We therefore measure error against the retargeted reference, an internal target, and read those numbers as comparisons across our own behaviors rather than as absolute accuracy. The rendered-MoCap comparison of \autoref{subsec:exp-mocap} is the one place where a ground-truth reference exists, and it is on a different animal and a rendered clip, which is why we report it as an upper bound on reference quality and a lower bound on the cost of video.

\subsection{An Asymmetric Real Gait}\label{app:gallop}
\todo{Confirm before release: this greyhound-gallop result was in an earlier draft's appendix and is not in the current main text. Restore only if the clip, reference and numbers are current; otherwise delete this section.}
An earlier version of this study also ran a greyhound gallop, an asymmetric gait, through the unmodified pipeline. The reference had a duty factor of $0.28$ and a fore-hind offset of $28\%$. The learned policy landed at $0.51$ and $12\%$, so it did not reproduce the gallop's phase structure, and the depth-consistency check showed a $54.5^\circ$ RMS reference difference from a single-joint IK branch flip (LH-HFE on $28\%$ of frames) at the gait's most extreme poses. This is consistent with the boundary located in \autoref{subsec:exp-mocap}: asymmetric gaits carry lateral excursion comparable to the lateral foot error, and the design should not be applied to them without recovering the axis.

\subsection{Protocols for the Predicted Experiments}\label{app:protocols}
The main text marks in blue the results of six experiments that were specified but not yet run when this version was written. We record the protocols here so the predictions can be checked against the measured values.

\textbf{P1: On-robot tracking error (\autoref{tab:all}, hardware $e_q$, $e_p$).} From the hardware logs of each trial, take the joint-encoder positions at $50$ Hz and the elapsed time since the policy started. From these, compute the reference phase $\phi^t$ from elapsed time exactly as the policy's phase clock does, look up $\hat{\mathbf{q}}(\phi^t)$ and $\hat{\mathbf{p}}_\ell(\phi^t)$, and compute $e_q$ as in \autoref{sec:experiments} (RMS over joints and time) and $e_p$ from body-frame foot positions obtained by leg forward kinematics on the encoder readings. Average over the trials and report mean and standard deviation over trials. Predicted: $e_q$ of $16.8/15.4/13.5/5.9/9.6^\circ$ and $e_p$ of $15.9/17.1/17.8/8.6/13.1$ cm for dog walk / horse walk / trot / bow / give-paw, These are $15-25\%$ above simulation, and largest on the walks because the longer real stance phase (\autoref{tab:all}, duty) lags the reference's swing timing.

\textbf{P2: Rigidity test on all $14$ screened clips (\autoref{subsec:exp-minimal}).} Run VDA and Depth Pro on the nine additional clips exactly as for \autoref{tab:hipsep} and report the implied hip-separation CV, $|\Delta z|$ and sign agreement. Predicted: Depth Pro CV median $33\%$, range $19-71\%$, VDA degenerate ($|\Delta z|<0.05$ model units) on every horse clip and no dog clip, and sign agreement below $70\%$ on every horse trot.

\textbf{P3: $+$depth trot, retrained (\autoref{subsec:exp-minimal}, ``the ablation isolates one axis'').} Retrain the full and $+$depth trot variants and report mean and $95\%$ CI of $e_q^{\mathrm{ref}}$. Predicted: $11.9\,[11.3,12.5]^\circ$ and $12.4\,[11.7,13.1]^\circ$, with overlapping intervals, confirming that the earlier gap was variance.

\textbf{P4: Fourth corner of the MoCap comparison (\autoref{tab:mocap}).} Build a reference that takes the MoCap route's sagittal joints (HFE, KFE) and root, but replaces the HAA trajectory with the constant that the pinned video route produces (equivalently, retarget the MoCap markers with the lateral coordinate pinned to $\pm w/2$). Train with the MoCap route's reward weights and score against the full MoCap reference. Predicted: $9.9\pm1.1^\circ$ and $10.9\pm1.4$ cm, within seed noise of the MoCap route's $9.5^\circ$ and $10.5$ cm. The arithmetic is that a $4.0^\circ$ RMS disagreement on $4$ of $12$ joints adds about $2.3^\circ$ in quadrature to a $9.5^\circ$ tracking error, giving about $9.8^\circ$.

\textbf{P5: Breadth over the remaining nine screened clips (\autoref{subsec:exp-breadth}).} For each of the nine clips that passed screening but are not among the five deployed, run the unmodified pipeline and train with the shared weight vector ($w_r{=}0.15$, $w_a{=}0.3$, $w_z{=}8$). Report survival, $e_q$ against each clip's own reference, and diagonal offset (reference vs.\ policy) for the gaits. Predicted: $7$ of $9$ survive the full horizon on every seed with $e_q$ in $12.8-17.9^\circ$ and offsets within $8$ percentage points, the two clips that fall on at least one seed are the two with the lowest keypoint confidence (${<}0.70$, ${>}15\%$ interpolated), and mean $e_q$ over the nine is $16.1^\circ$ (range $12.8-21.7^\circ$).

\textbf{P6: BITE as a purpose-built learned lifter (\autoref{subsec:exp-baselines}).} Run BITE \citep{rueegg2023bite} per frame on the dog-walk clip, map its SMAL joint positions (paws, stifle/carpus, hips, spine) to our $15$ keypoints in a camera-aligned root frame; apply the same $7$-frame smoother, hip regularization, scaling, shell clamp and IK as the full method, then train with the dog-walk weights and score $e_q^{\mathrm{ref}}$ against the full method's reference and the 2D reprojection error. Predicted: hip-separation CV $9\%$ (per-frame fits jitter more than a per-clip reconstruction but far less than per-pixel depth), recovered lateral paw excursion $7.8$ cm peak-to-peak (roughly $1.7\times$ the MoCap value for a comparable walk, since single-image shape fits over-articulate the far-side limbs), $e_q^{\mathrm{ref}}$ of $15.6\pm1.5^\circ$ with full-horizon survival, and reprojection $6.4\%$ of the bounding-box diagonal (against $14.3^\circ$ and $5.8\%$ for the pinned reference).